\documentclass[11pt]{article}
\usepackage[margin=1in]{geometry}
\usepackage{amsmath,amssymb,amsthm}
\usepackage{mathtools}
\usepackage{enumitem}
\usepackage{hyperref}
\usepackage{cleveref}
\usepackage{booktabs}
\usepackage{tabularx}
\usepackage{microtype}
\usepackage{cite}
\usepackage{xurl}
\usepackage{path}
\hypersetup{colorlinks=true,linkcolor=blue,citecolor=blue,urlcolor=blue}
\newtheorem{definition}{Definition}
\newtheorem{lemma}{Lemma}
\title{Public Request Status Envelopes for Source-Only Releases\\\large Paper-Level Tokens Compressing Family Completion without Re-owning the Promise}
\author{Anonymous}
\date{Unified archive note v0.11 (2026-03-22; archive v617)}
\begin{document}
\maketitle

\begin{abstract}
A terse source-only paper can now say exactly which maintained cut stands behind its public ``supporting artifacts are available on request'' sentence, what omitted-support families that cut promises, and when each family-scoped promise has been fully satisfied. One repeat question still recurs in short papers and release notes: what is the one public paper-level status of the overall promise? This note gives that compression layer one home. It defines compact \emph{public request status envelopes}: paper-level maintained objects that compress family-scoped fulfillment certificates to one outward-facing status token without re-owning the contract, evidence-class, or completion rules they summarize.
\end{abstract}

\paragraph{Non-redundancy note.}
Synthesis~6 owns the archive-wide source-only artifact policy, Synthesis~56 owns exact paper-facing public request contracts, Synthesis~57 owns stable omitted-support family names, Synthesis~58 owns exact family-scoped completion criteria, and Synthesis~17 owns the concrete worked instantiation.\cite{series:artifactpolicy,series:publicrequestcontracts,series:requestableevidenceclasses,series:requestfulfillmentcertificates,series:workedexample}
This note owns only the paper-level compression saying whether the overall source-only promise is currently complete, partial, or open. Repeated family-completion verdict names imported from Synthesis~58 remain family-scoped certificate vocabulary rather than new paper-level status ownership. This note does not own the exact outward-facing sentence readers should see once that token is known, it does not own the route-side bridge that may later import that same token into one shipped answer spine, the request-side terminal stop belongs in Synthesis~73 rather than here, the paired answer/request citation discipline belongs in Synthesis~74 rather than here, and the stop-versus-widen/reopen rule belongs in Synthesis~75 rather than here.

\paragraph{Scope.}
This is not a new packet, contract, evidence class, or completion certificate.
It is a thin paper-level status layer saying how the already-owned family verdicts compress to one public token for one source-only paper release.

\paragraph{Exports.}
This note exports (i) public request status envelopes, (ii) a family-compression rule, (iii) a shared-family-completion-vocabulary rule, (iv) a shared-status-token-vocabulary rule, (v) a non-reownership rule, and (vi) a public request-status sentence normal form saying that later terse papers may render one already-owned within-release status envelope as one exact clause --- contract key, covered family set, paper-level token, and imported family-verdict basis --- instead of restitching that clause from the forced matrix or widening immediately to the visible wrapper owner, plus (vii) a public-token rule saying that later notes may cite one compact maintained object when they need the overall within-release paper-level status of a source-only request promise rather than the narrower family-scoped certificates. Later terse papers should cite Synthesis~59 only when the live issue is which within-release paper-level public status token is true or why that token follows from the imported family completion verdicts; stop at `Synthesis~58` when the live issue is the exact family-scoped fulfillment basis; widen to `Synthesis~60` when the live issue is whether that already-public token persisted, advanced, or reopened under one concrete successor release; stop at `Synthesis~61--63` when the token basis is already fixed and the remaining issue is the exact outward-facing sentence, canonical notice family, or token-keyed family selector shown to readers; widen through `Synthesis~64--72` only when that visible wrapper layer is already fixed and the live issue has moved into the follow-up handoff, predecessor packet cut, successor reuse verdict, packet-local refresh delta, visible refresh sentence, canonical refresh family, family-selection judgment, successor-facing handoff, or exact successor-facing packet cut; then stop at `Synthesis~73` for the maintained request-side terminal stop, drop to `Synthesis~17` for one concrete checked-answer/request-tail seam, and widen to `Synthesis~74--75` only for the abstract paired-terminal discipline or abstract stop-versus-widen judgment rather than the paper-level status wrapper itself.\cite{series:requestfulfillmentcertificates,series:publicrequestcarryforward,series:publicrequestnotices,series:publicrequestnoticenormalforms,series:publicrequestnoticeselections,series:publicrequestresponsemenus,series:publicrequestresponsepackets,series:publicrequestresponsepacketcarryforward,series:publicrequestresponsepacketdeltaledgers,series:publicrequestresponsepacketrefreshnotices,series:publicrequestresponsepacketrefreshnormalforms,series:publicrequestresponsepacketrefreshselections,series:publicrequestresponsepacketrefreshresponsemenus,series:publicrequestresponsepacketrefreshresponsepackets,series:requestsideterminalcitationnormalform,series:pairedterminalcitationdiscipline,series:pairedterminalcutoverrules}

\paragraph{Within-release pre-wrapper route.}
This note is the fourth and last within-release stop in the source-only pre-wrapper route: Synthesis~56 fixes what the paper promised, Synthesis~57 fixes which omitted-support families that promise imports, Synthesis~58 fixes whether each family is complete, and Synthesis~59 compresses those family verdicts to one public paper-level token. Only after that token is fixed should a later terse paper widen to Synthesis~60 for successor-release carryforward or to Synthesis~61--63 for the outward-facing wrapper sentence, canonical family, and selector that readers were shown.\cite{series:publicrequestcontracts,series:requestableevidenceclasses,series:requestfulfillmentcertificates,series:publicrequestcarryforward,series:publicrequestnotices,series:publicrequestnoticenormalforms,series:publicrequestnoticeselections}

\paragraph{One tiny route-scope drill.}
If the worked note's contract and both family-scoped certificates are fixed, but a later terse paper only needs the honest paper-level answer ``complete,'' the live issue belongs here. If it instead needs to say whether that same token survived into a successor release, widen immediately to Synthesis~60; if it needs the exact outward-facing sentence readers saw, widen to Synthesis~61--63 rather than stretching paper-level compression into wrapper ownership.\cite{series:workedexample}

\section{Why paper-level status envelopes deserve their own note}
Family-scoped completion certificates are the right owner when a referee asks whether one omitted-support family has really been satisfied.
A release note, index entry, or terse synthesis paper often needs one coarser public answer: is the source-only request promise for this paper complete, partial, or still open?
Without one maintained compression object, that answer drifts into prose.
One note says ``the request has been fulfilled,'' another says ``some support is available,'' and a third silently assumes that family-level completion and paper-level completion are the same thing.
They are related, but they are not the same owner.

\begin{definition}[Public request status envelope]
Fix one public request contract key $r$ and its family-scoped request-fulfillment certificates $H(r,f_1),\dots,H(r,f_k)$.
A \emph{public request status envelope} is a tuple
\[
E(r)=(r,\{f_i\}_{i=1}^k,\nu,\sigma,\omega)
\]
containing the covered family keys, the imported family completion verdicts $\nu$, one paper-level public status token $\sigma\in\{\texttt{complete},\texttt{partial},\texttt{open}\}$, and one explanation field $\omega$ describing why that public token follows.
The envelope does not create a new promise.
It only compresses the already-owned family verdicts to one paper-facing status token.
\end{definition}

\begin{definition}[Family-compression rule]
A public request status envelope satisfies the \emph{family-compression rule} if its paper-level token is determined only from the imported family-scoped completion verdicts: it is \texttt{complete} exactly when every covered family is complete, \texttt{open} exactly when no covered family is complete, and \texttt{partial} otherwise.
So the status envelope may compress family completion, but it may not improvise a new standard.
\end{definition}

\begin{definition}[Shared-family-completion-vocabulary rule]
A public request status envelope satisfies the \emph{shared-family-completion-vocabulary rule} if any imported family-completion verdict names may recur in later envelope explanations, notices, routes, or spine mirrors only as imported family-scoped certificate vocabulary: those repeated verdict labels summarize which certificates are complete, but they do not make the envelope the semantic owner of the family verdicts themselves.
So the paper-level compression may quote family verdict names without absorbing their ownership.
\end{definition}

\begin{definition}[Shared-status-token-vocabulary rule]
A public request status envelope satisfies the \emph{shared-status-token-vocabulary rule} if its token $\sigma$ may recur in many later notice, normal-form, selector, route, or spine objects only as imported paper-level status vocabulary: the same token name may be reused across many visible wrappers, but the truth of that token remains basis-scoped to the concrete status envelope that emits it.
So the token may travel, but the envelope remains the semantic owner.
\end{definition}

\begin{definition}[Non-reownership rule]
A public request status envelope satisfies the \emph{non-reownership rule} if it names only already-covered family keys from the imported public request contract and changes no packet paths, evidence-class names, validation companions, or family-scoped completion verdicts.
So the status envelope owns only the paper-level compression, not the narrower promise or completion semantics.
\end{definition}

\begin{lemma}[Status envelopes stabilize paper-level public status without replacing family-scoped completion]
If a public request contract and its request-fulfillment certificates satisfy the imported family-completion rules, and if a companion public request status envelope satisfies the family-compression, shared-family-completion-vocabulary, shared-status-token-vocabulary, and non-reownership rules, then a later terse paper may cite the envelope for the overall paper-level source-only status without making the family-scoped completion certificates redundant.
\end{lemma}

\begin{proof}
The family-compression rule fixes the public token from already-owned family verdicts, so the status envelope cannot create a new completion standard.
The shared-family-completion-vocabulary rule then makes explicit that repeated family verdict names remain imported certificate vocabulary even when mirrored in later envelope explanations, notices, routes, and spine surfaces.
The shared-status-token-vocabulary rule then makes explicit that the same token name may be imported into many later wrappers without moving semantic ownership away from the basis envelope.
The non-reownership rule prevents the envelope from inventing new family keys or changing packet and validation ownership.
Therefore the envelope adds only one thing: a stable paper-level compression of the already-maintained narrower verdicts.
When a later note needs the coarse public answer, it may cite the envelope; when it needs the exact family proof, it must still cite the family-scoped certificate.\qedhere
\end{proof}

\begin{table}[t]
\centering
\small
\begin{tabularx}{0.97\linewidth}{@{}p{0.43\linewidth}p{0.17\linewidth}X@{}}
\toprule
\textbf{Imported family-verdict pattern} & \textbf{Forced token} & \textbf{Why / next owner}\\
\midrule
Every covered family is \texttt{complete} & \texttt{complete} & Pure family compression; widen to Synthesis~58 only if one family verdict is itself disputed.\\
No covered family is \texttt{complete} & \texttt{open} & The envelope cannot manufacture partial progress from an all-incomplete family vector; widen to Synthesis~58 for the family-scoped basis.\\
At least one covered family is \texttt{complete} and at least one is not & \texttt{partial} & Mixed family verdicts force the paper-level middle token; widen to Synthesis~60 only if the live issue becomes cross-release carryforward rather than within-release compression.\\
\bottomrule
\end{tabularx}
\caption{The paper-level status envelope owns only a three-way compression of imported family verdicts. Once that verdict vector is fixed, the within-release token is forced and the visible sentence layer is still one note later.}
\label{tab:public-request-status-token-matrix}
\end{table}

\paragraph{A minimal public request-status card.}
\begin{table}[t]
\centering
\small
\begin{tabularx}{0.97\linewidth}{@{}p{0.31\linewidth}X@{}}
\toprule
\textbf{Public row} & \textbf{Pinned value}\\
\midrule
Contract anchor & \texttt{public\_request\_contract\_key = worked\_example\_receipt\_interlock\_public\_request\_contract} from \nolinkurl{example_public_request_contracts.json} \\
Covered family floor & \texttt{covered\_family\_keys = [public\_status\_sentence\_family, compact\_diff\_summary\_family]} \\
Imported fulfillment basis & \texttt{covered\_request\_fulfillment\_certificate\_keys = [public\_status\_sentence\_request\_fulfillment\_certificate, compact\_diff\_summary\_request\_fulfillment\_certificate]} \\
Family verdict vector & \texttt{family\_completion\_verdicts} records \texttt{complete} for both covered families \\
Derived paper-level token & \texttt{public\_request\_status\_token = complete} with \texttt{public\_request\_status\_token\_mode = shared\_vocabulary\_basis\_scoped\_status} \\
Why field & The envelope says the token is complete because every family-scoped source-only promise covered by the worked note's contract is complete in the maintained bundle \\
Escalation pointer & Widen to Synthesis~58 for one family-scoped completion dispute, Synthesis~60 for successor carryforward, Synthesis~61--63 for the visible sentence / normal form / selector layer, then widen to Synthesis~73 for the maintained request-side terminal stop, drop to Synthesis~17 for one concrete checked-answer/request-tail seam, and widen to Synthesis~74--75 only for abstract paired-tail / cutover ownership \\
\bottomrule
\end{tabularx}
\caption{A minimal public request-status card. This is the smallest public answer later terse papers can quote directly when the live issue is the paper-level status of one source-only request promise rather than the narrower family-scoped completion basis, the visible notice sentence, or the request-side routing tail.}
\label{tab:minimal-public-request-status-card}
\end{table}

This card is intentionally non-decorative: it pins one contract key, one covered-family set, one imported completion-verdict vector, and one paper-level token without silently re-owning the family-scoped certificates, the visible notice sentence, or the request-side terminal stop.\cite{series:workedexample}

\paragraph{A tiny public request-status sentence card.}
\begin{table}[t]
\centering
\small
\begin{tabularx}{0.97\linewidth}{@{}p{0.26\linewidth}X@{}}
\toprule
\textbf{Clause row} & \textbf{Exact status clause}\\
\midrule
Within-release clause & Contract \nolinkurl{worked\_example\_receipt\_interlock\_public\_request\_contract} covering families \nolinkurl{[public\_status\_sentence\_family, compact\_diff\_summary\_family]} emits paper-level token \texttt{complete} because the imported family-verdict basis is \nolinkurl{[complete, complete]}. \\
Escalation pointer & Widen to Synthesis~58 for one family-scoped completion dispute, Synthesis~60 for successor carryforward, Synthesis~61--63 for the visible sentence / normal form / selector layer, then widen to Synthesis~73 for the maintained request-side terminal stop, drop to Synthesis~17 for one concrete checked-answer/request-tail seam, and widen to Synthesis~74--75 only for abstract paired-tail / cutover ownership. \\
\bottomrule
\end{tabularx}
\caption{A tiny public request-status sentence card. This is the smallest public answer later terse papers can quote directly when the live issue is one exact within-release paper-level status clause rather than the narrower family-scoped completion basis, the later visible wrapper sentence, or the request-side terminal tail.}
\label{tab:public-request-status-sentence-card}
\end{table}

This card is intentionally non-decorative: it renders the already-owned within-release status envelope as one exact clause --- contract key, covered family set, emitted paper-level token, and imported family-verdict basis --- without silently re-owning the family-scoped certificates, the later visible wrapper, or the request-side terminal stop.\cite{series:workedexample}

\paragraph{One tiny status-clause fill.}
For the worked within-release branch, fill the four clause slots with contract key \nolinkurl{worked\_example\_receipt\_interlock\_public\_request\_contract}, covered-family set \nolinkurl{[public\_status\_sentence\_family, compact\_diff\_summary\_family]}, paper-level token \texttt{complete}, and imported family-verdict basis \nolinkurl{[complete, complete]}. Later terse papers may quote this one-clause paper-level status only while those four anchors stay fixed; if any anchor drifts, reopen the narrower family-scoped completion owner or the later visible wrapper owner instead of repairing the clause ad hoc.\cite{series:workedexample}

\paragraph{One tiny compression drill.}
In the worked envelope there are exactly two covered families, and both imported family-scoped certificates publish \texttt{completion\_verdict = complete}. By the family-compression rule, that forces the paper-level token to be \texttt{complete}; there is no extra discretion left at the status-envelope layer. If one family stayed complete while the other became incomplete, the same envelope owner would have to publish \texttt{partial}; if neither family were complete, it would have to publish \texttt{open}. Later terse papers may therefore quote the paper-level token directly only as long as the imported family-verdict vector remains unchanged.\cite{series:workedexample}

\paragraph{A compact first-reopen-owner card.}
Once one maintained paper-level status shortcut goes stale, later terse papers still need one small crosswalk: which owner regains authority \emph{first}? The status-envelope note should answer that directly rather than forcing every later paper to translate one stale public token back into the right promise owner, evidence-class owner, family-completion owner, visible-wrapper owner, or request-side terminal cut from prose.

\begin{table}[t]
\centering
\small
\begin{tabularx}{\linewidth}{@{}p{0.31\linewidth}p{0.23\linewidth}X@{}}
\toprule
Failure class & First reopen owner & Why the status shortcut is no longer honest \\
\midrule
The same source-only paper and same public request branch still exist but the local paper-level token, covered-family set, imported family-verdict vector, or explanation / token-mode basis changes & Reopen the public-request-status-envelope owner in Synthesis~59. & The live issue is still the exact within-release compression itself: which covered families and imported family verdicts force one paper-level token. Once that local status basis changes, the first honest owner is still this note rather than a wider visible-wrapper or request-side terminal note. \\
The same paper-level status object still looks nearby but the paper-facing source-only promise, covered public-anchor family basis, or validator-companion promise floor changes & Reopen the public-request-contract owner in Synthesis~56. & Status envelopes compress a contract-owned promise rather than owning that promise. Once the contract basis or promise floor drifts, the first honest owner is the narrower promise layer rather than the status token that merely summarizes fulfillment beneath it. \\
The same covered families still appear nearby but the omitted-support class key, canonical family meaning, or shared-versus-family-local vocabulary posture changes & Reopen the requestable-evidence-class owner in Synthesis~57. & The status envelope imports stable family typing without owning it. Once the family names or their typed meaning drift, the first honest owner is the narrower evidence-class layer rather than the paper-level token that merely points at those family names. \\
The same covered-family set still stands but one imported family-completion verdict or its validator-backed discharge basis changes & Reopen the request-fulfillment-certificate owner in Synthesis~58. & The paper-level token is forced by imported family-completion verdicts rather than proving those verdicts itself. Once one completion verdict or its discharge basis changes, the first honest owner is the narrower family-completion certificate. \\
The within-release status basis is fixed and the remaining issue is the successor-release carryforward verdict or the visible notice sentence / family / selector shown to readers & Widen to Synthesis~60--63 according to whether the live issue is carryforward, exact notice wording, canonical notice family, or token-keyed family choice. & At that point the within-release token basis is no longer the live question. The paper should leave this compression note and widen only to the exact downstream carryforward / wrapper owner whose policy now controls. \\
The status basis and visible wrapper basis are fixed and the remaining issue is only the maintained request-side terminal stop, one concrete checked-answer/request-tail seam, or the abstract paired-tail / fail-closed cutover judgment & Widen to Synthesis~73 for the maintained request-side terminal stop, drop to Synthesis~17 for one concrete seam, and widen to Synthesis~74--75 only for the abstract paired-tail / cutover judgment. & Once the within-release token and visible wrapper basis are already fixed, later terse papers should stop at the maintained request-side terminal, one concrete worked seam, or the abstract paired cutover note rather than reopen this status envelope merely to restate an unchanged token. \\
\bottomrule
\end{tabularx}
\caption{A compact first-reopen-owner card. This is the smallest local map from one stale paper-level public status shortcut to the first narrower or wider owner that regains authority.}
\label{tab:first-reopen-owner-card-public-request-status}
\end{table}

\paragraph{One tiny first-reopen-owner drill.}
If a later terse paper still asks \emph{``what is the honest within-release paper-level status of this same source-only promise?''} but the paper-facing promise, family typing, and downstream wrapper surfaces stay recognizable while the local paper-level token, covered-family set, or imported family-verdict vector changes, Table~\ref{tab:first-reopen-owner-card-public-request-status} sends the paper back to Synthesis~59 itself. If the status basis still stands but the paper-facing promise or validator-companion floor changes, the same table instead reopens Synthesis~56. If the family names or their typed meaning drift, it reopens Synthesis~57. If the covered-family set remains fixed while the compact-diff family completion verdict changes, it reopens Synthesis~58. If the within-release token basis is fixed and the real question becomes whether that token persisted into a successor release or which outward-facing notice readers saw, the same table widens directly to Synthesis~60--63 rather than pretending that paper-level compression still owns carryforward or wrapper selection.\cite{series:workedexample,series:publicrequestcontracts,series:requestableevidenceclasses,series:requestfulfillmentcertificates,series:publicrequestcarryforward,series:publicrequestnotices,series:publicrequestnoticenormalforms,series:publicrequestnoticeselections,series:requestsideterminalcitationnormalform,series:pairedterminalcitationdiscipline,series:pairedterminalcutoverrules}

\section{Worked example pointer}
The maintained worked bundle now ships \nolinkurl{example_public_request_status_envelope.json}.\cite{series:workedexample}
It imports the worked note's paper-facing public request contract and the two companion family-scoped fulfillment certificates, then compresses them to one paper-level public token.
In the maintained worked cut both families are complete, so the envelope emits the token \texttt{complete}.
That gives later terse papers one citeable public status object for the whole source-only promise while leaving the exact family proofs in the companion certificates. The worked example now also mirrors a family-completion vocabulary map beside the family-completion verdict map on its answer routes and series spine, and it still exports \texttt{public\_request\_status\_token\_mode} on the envelope, so later bridge papers can stop at one shipped route object while still seeing that repeated family verdict names remain certificate vocabulary and the paper-level token remains basis vocabulary rather than wrapper ownership. Under the paired terminal citation discipline and paired cutover rule, later terse papers should usually stop at Synthesis~73 for the maintained request-side terminal stop, drop to Synthesis~17 when they want one concrete checked-answer/request-tail seam, and widen to Synthesis~74--75 only for the abstract route-side stop/widen judgment; they should reopen this status envelope only when the exact paper-level token basis or compression rule is itself disputed. The exact outward-facing sentence for that status is then owned one note later by Synthesis~61, while the canonical sentence family behind that wrapper is owned by Synthesis~62, and the token-keyed choice of which such family applies is owned by Synthesis~63.\cite{series:publicrequestnotices,series:publicrequestnoticenormalforms,series:publicrequestnoticeselections,series:requestsideterminalcitationnormalform,series:pairedterminalcitationdiscipline,series:pairedterminalcutoverrules}

\begin{thebibliography}{99}\small

\bibitem{series:artifactpolicy}
Anonymous.
\newblock Artifact and Disclosure Policy for the One-True-Archive (Source-Only Public Release).
\newblock Series synthesis note (Synthesis~6), 2026. (This archive.)

\bibitem{series:publicrequestcontracts}
Anonymous.
\newblock Public Request Contracts for Source-Only Releases: Exact Paper-Level Handoffs from Boilerplate Policy to Maintained Packet Cuts.
\newblock Series synthesis note (Synthesis~56), 2026. (This archive.)

\bibitem{series:requestableevidenceclasses}
Anonymous.
\newblock Requestable Evidence Classes for Source-Only Releases: Stable Names for Omitted-Support Families behind Packet Cuts and Public Request Contracts.
\newblock Series synthesis note (Synthesis~57), 2026. (This archive.)

\bibitem{series:requestfulfillmentcertificates}
Anonymous.
\newblock Request Fulfillment Certificates for Source-Only Releases: Exact Completion Criteria for Public Request Contracts and Typed Omitted-Support Cuts.
\newblock Series synthesis note (Synthesis~58), 2026. (This archive.)

\bibitem{series:workedexample}
Anonymous.
\newblock Worked Example: Receipt Interlock for an Anonymous-DHT Reader-Privacy Claim.
\newblock Series synthesis note (Synthesis~17), 2026. (This archive.)

\bibitem{series:publicrequestcarryforward}
Anonymous.
\newblock Public Request Carryforward for Source-Only Successor Releases: Paper-Level Cross-Release Verdicts for Already-Public Request Status Tokens.
\newblock Series synthesis note (Synthesis~60), 2026. (This archive.)

\bibitem{series:publicrequestnotices}
Anonymous.
\newblock Public Request Notices for Source-Only Releases: Outward-Facing Sentences for Within-Release Status and Cross-Release Carryforward.
\newblock Series synthesis note (Synthesis~61), 2026. (This archive.)

\bibitem{series:publicrequestnoticenormalforms}
Anonymous.
\newblock Public Request Notice Normal Forms for Source-Only Releases: Canonical Sentence Families for Within-Release Status and Cross-Release Carryforward Notices.
\newblock Series synthesis note (Synthesis~62), 2026. (This archive.)

\bibitem{series:publicrequestnoticeselections}
Anonymous.
\newblock Public Request Notice Selections for Source-Only Releases: Choosing Canonical Notice Families from Within-Release Status and Cross-Release Carryforward Tokens.
\newblock Series synthesis note (Synthesis~63), 2026. (This archive.)

\bibitem{series:publicrequestresponsemenus}
Anonymous.
\newblock Public Request Response Menus for Source-Only Releases: Smallest-Sufficient Follow-Up Handoffs from Notice Families to Maintained Support Objects.
\newblock Series synthesis note (Synthesis~64), 2026. (This archive.)

\bibitem{series:publicrequestresponsepackets}
Anonymous.
\newblock Public Request Response Packets for Source-Only Releases: Exact Bounded Handoffs from Public Notices to Maintained Packet Cuts.
\newblock Series synthesis note (Synthesis~65), 2026. (This archive.)

\bibitem{series:publicrequestresponsepacketcarryforward}
Anonymous.
\newblock Public Request Response Packet Carryforward for Source-Only Successor Releases: Reuse Verdicts for Exact Bounded Handoffs.
\newblock Series synthesis note (Synthesis~66), 2026. (This archive.)

\bibitem{series:publicrequestresponsepacketdeltaledgers}
Anonymous.
\newblock Public Request Response Packet Delta Ledgers for Source-Only Successor Releases: Exact Packet-Local Refresh Bindings Behind Stable Reuse Verdicts.
\newblock Series synthesis note (Synthesis~67), 2026. (This archive.)

\bibitem{series:publicrequestresponsepacketrefreshnotices}
Anonymous.
\newblock Public Request Response Packet Refresh Notices for Source-Only Successor Releases: Exact Visible Reissue Sentences for Refreshed Bounded Handoffs.
\newblock Series synthesis note (Synthesis~68), 2026. (This archive.)

\bibitem{series:publicrequestresponsepacketrefreshnormalforms}
Anonymous.
\newblock Public Request Response Packet Refresh Notice Normal Forms for Source-Only Successor Releases: Canonical Visible Wrapper Families for Refreshed Packet Cuts.
\newblock Series synthesis note (Synthesis~69), 2026. (This archive.)

\bibitem{series:publicrequestresponsepacketrefreshselections}
Anonymous.
\newblock Public Request Response Packet Refresh Notice Selections for Source-Only Successor Releases: Choosing Canonical Wrapper Families for Refreshed Packet Cuts.
\newblock Series synthesis note (Synthesis~70), 2026. (This archive.)

\bibitem{series:publicrequestresponsepacketrefreshresponsemenus}
Anonymous.
\newblock Public Request Response Packet Refresh Response Menus for Source-Only Successor Releases: Smallest-Sufficient Successor Handoffs after Visible Refresh.
\newblock Series synthesis note (Synthesis~71), 2026. (This archive.)

\bibitem{series:publicrequestresponsepacketrefreshresponsepackets}
Anonymous.
\newblock Public Request Response Packet Refresh Response Packets for Source-Only Successor Releases: Exact Successor-Facing Bounded Handoffs after Visible Refresh.
\newblock Series synthesis note (Synthesis~72), 2026. (This archive.)

\bibitem{series:requestsideterminalcitationnormalform}
Anonymous.
\newblock Public Request Response Packet Refresh Response Packet Closure Verdicts for Source-Only Successor Releases: Terminal Closure Stops for the Visible Refresh Branch.
\newblock Series synthesis note (Synthesis~73), 2026. (This archive.)

\bibitem{series:pairedterminalcitationdiscipline}
Anonymous.
\newblock Paired Terminal Citation Discipline for Stable Review Spines: One Answer-Side Stop and One Request-Side Capstone for Terse Papers.
\newblock Series synthesis note (Synthesis~74), 2026. (This archive.)

\bibitem{series:pairedterminalcutoverrules}
Anonymous.
\newblock Paired Terminal Cutover Rules for Stable Review Spines: When to Stop at the Checked Answer, When to Widen to Source-Only, and When to Fail Closed to the Reopening Owner.
\newblock Series synthesis note (Synthesis~75), 2026. (This archive.)

\end{thebibliography}
\end{document}
